\usepackage[T1]{fontenc} %pacote inicial fonte
\usepackage[margin=1in]{geometry} % para margem principalmente, mas tambem formatar texto
\usepackage{xcolor} % para cor de texto
\usepackage{lipsum} % para texto aleatorio
\usepackage[latin,brazilian]{babel} % para linguas entendidas pelo texto
\usepackage{graphicx} % para imagens
\usepackage{array} %para tabelas
\usepackage{booktabs} %para tabelas tambem, so que mais completo
\usepackage{amsmath} % para matematica
\usepackage{mathcomp}
\usepackage{amsfonts} % para formatacao de texto matematico
\usepackage{amssymb}
\usepackage{url}
\usepackage{hyperref}
\usepackage{setspace}
\usepackage[parfill]{parskip} % para espaçamento no texto entre paragrafos
\usepackage{listings} % para trechos de codigo PySpark/Python
\usepackage{float} % para posicionamento [H] de figuras
\usepackage{natbib} % para citacoes tambem exige que os arquivos bib sejam excluidos
%\usepackage[style=authoryear]{biblatex} % tomar cuidado que biblatex não tem suporte para o latin e pode exigir recompilacao do bib
\usepackage{csquotes}
%\addbibresource{learnlatex.bib} %file of reference

% Configuracao de estilo para listagens de codigo PySpark
\definecolor{codeverde}{rgb}{0.0,0.5,0.0}
\definecolor{codecinza}{rgb}{0.5,0.5,0.5}
\definecolor{codelaranja}{rgb}{0.7,0.35,0.0}
\definecolor{codefundo}{rgb}{0.97,0.97,0.97}
\lstset{
  language=Python,
  basicstyle=\ttfamily\footnotesize,
  keywordstyle=\color{blue}\bfseries,
  commentstyle=\color{codeverde}\itshape,
  stringstyle=\color{codelaranja},
  numbers=left,
  numberstyle=\tiny\color{codecinza},
  numbersep=8pt,
  stepnumber=1,
  frame=single,
  rulecolor=\color{codecinza},
  backgroundcolor=\color{codefundo},
  breaklines=true,
  breakatwhitespace=true,
  showstringspaces=false,
  tabsize=4,
  captionpos=b,
  extendedchars=true,
  literate=%
    {á}{{\'a}}1 {é}{{\'e}}1 {í}{{\'i}}1 {ó}{{\'o}}1 {ú}{{\'u}}1
    {Á}{{\'A}}1 {É}{{\'E}}1 {Í}{{\'I}}1 {Ó}{{\'O}}1 {Ú}{{\'U}}1
    {â}{{\^a}}1 {ê}{{\^e}}1 {ô}{{\^o}}1 {Â}{{\^A}}1 {Ê}{{\^E}}1 {Ô}{{\^O}}1
    {ã}{{\~a}}1 {õ}{{\~o}}1 {Ã}{{\~A}}1 {Õ}{{\~O}}1
    {à}{{\`a}}1 {ç}{{\c c}}1 {Ç}{{\c C}}1
}
